\documentclass{article}
\usepackage{amsmath}
\usepackage{amssymb}
\usepackage{geometry}
\geometry{margin=1in}

\title{A Novel $\phi$-Based Numeral System Derived from 6-Basis Vectors}
\author{James A. Hinds}
\date{\today}

\begin{document}

\maketitle

\begin{abstract}
We introduce a numeral system based on a 6-basis vector representation, in which each component represents an exponent of the golden ratio ($\phi$). Unlike traditional positional numeral systems, increments are defined by integer powers of $\phi$, which satisfies $\phi^n = \phi^{n-1} + \phi^{n-2}$, with $\phi^2 = \phi + 1$. Each valid 6-vector corresponds uniquely to an intersection point in Cartesian space. This system has potential applications in computational mathematics, geometric modeling, and theoretical exploration.
\end{abstract}

\section{Introduction}
Traditional numeral systems rely on integer bases such as 2, 10, or 16. In contrast, we propose a numeral system based on the golden ratio $\phi$, an irrational number with unique algebraic properties. Numbers are represented via 6-basis vectors, each component encoding a power of $\phi$. A unique Cartesian intersection criterion defines the validity of these vectors.

\section{Theoretical Background}

\subsection{The Golden Ratio and Recurrence}
The golden ratio $\phi$ is defined by:
\[
\phi = \frac{1 + \sqrt{5}}{2},
\]
satisfying the recurrence:
\[
\phi^n = \phi^{n-1} + \phi^{n-2}, \quad \phi^2 = \phi + 1.
\]

\subsection{6-Basis Vectors and Intersection}
We define six basis vectors in $\mathbb{R}^3$:
\[
\text{sixBases} = \{[\phi,0,1], [\phi,0,-1], [0,1,\phi], [0,-1,\phi], [1,\phi,0], [-1,\phi,0]\}.
\]

A 6-vector $\mathbf{v} = [v_0, \dots, v_5]$ is valid if it yields a unique intersection point via:
\[
\text{sixToCartesian}(\mathbf{v}) = \sum_{i=0}^{5} v_i \cdot \text{sixBases}[i].
\]

We specifically consider partial vectors of the form $[a, ?, c, ?, e, ?]$ where unknown components satisfy:
\begin{align*}
b &= \frac{2e}{\phi} - a,\\
d &= -c,\\
f &= -e.
\end{align*}

Thus, the intersection coordinates become:
\[
X = 4e, \quad Y = 2c, \quad Z = 2\left(a - \frac{e}{\phi}\right).
\]

\section{Methodology}
Our procedure includes two key steps:
\begin{enumerate}
\item Completing partial vectors using:
\[ b = \frac{2e}{\phi} - a, \quad d = -c, \quad f = -e. \]

\item Mapping vectors to Cartesian points via:
\[
\text{sixToCartesian}(\mathbf{v}) = \sum_{i=0}^{5} v_i \cdot \text{sixBases}[i].
\]
The consistency is validated against expected coordinates.
\end{enumerate}

\section{Results and Discussion}
Preliminary computational tests demonstrate the robustness of this system. Visualization tools confirm the unique encoding of intersection points. The numeral representation introduces a non-integer radix based explicitly on $\phi$, offering a fresh perspective on numeral system structures.

\section{Conclusion}
We have developed a novel numeral system based on a 6-basis vector framework using $\phi$. The strict Cartesian intersection criterion guarantees unique number representation. Future research includes refining the arithmetic framework and exploring applications in computational mathematics, cryptography, and quantum computing.

\section*{Acknowledgments}
[Include acknowledgments here.]


\section*{Appendix: On Arithmetic in the $\phi$-Based System}

The arithmetic (addition and subtraction) rules for this numeral system remain an open challenge. Unlike standard numeral systems, arithmetic in this $\phi$-based system involves bidirectional carry propagation, as demonstrated by the example:
\[
2 \neq \phi^2 + \phi^{-1} \approx 3.236,
\]
but rather:
\[
2 = \phi + \phi^{-2} \approx 2.0.
\]
Future work must rigorously define arithmetic operations, normalization, and stability criteria to ensure the consistency and uniqueness of numeral representations.

i\vfill
\smallskip
\noindent\small © 2025 James A. Hinds . All rights reserved.

\end{document}

